% ============================================================
%  DATA STORM 7.0 — STORMING ROUND SUBMISSION REPORT
%  Team: BigBug
%  Max 5 pages including cover page
% ============================================================

\documentclass[11pt, a4paper]{article}

% ── Packages ────────────────────────────────────────────────
\usepackage[T1]{fontenc}        % Better font encoding
\usepackage{mathpazo}           % Palatino font for main text (elegant and professional)
\usepackage{inconsolata}        % Clean, modern monospace font for \texttt{}
\usepackage[
  top=2cm, bottom=2cm,
  left=2.2cm, right=2.2cm
]{geometry}
\usepackage{graphicx}
\usepackage{xcolor}
\usepackage{booktabs}
\usepackage{tabularx}
\usepackage{array}
\usepackage{enumitem}
\usepackage{titlesec}
\usepackage{fancyhdr}
\usepackage{hyperref}
\usepackage{amsmath}
\usepackage{amssymb}
\usepackage{multicol}
\usepackage{parskip}
\usepackage{microtype}
\usepackage{tcolorbox}
\usepackage{fontawesome5}

% ── Colours (customise to your team palette) ────────────────
\definecolor{primary}{HTML}{1A3C5E}
\definecolor{accent}{HTML}{E87722}
\definecolor{lightgray}{HTML}{F4F5F7}
\definecolor{midgray}{HTML}{6B7280}
\definecolor{successgreen}{HTML}{16A34A}
\definecolor{alertred}{HTML}{DC2626}

% ── Section formatting ───────────────────────────────────────
\titleformat{\section}
  {\large\bfseries\color{primary}}
  {\thesection.}{0.5em}{}
  [\color{accent}\titlerule[1.2pt]]

\titleformat{\subsection}
  {\normalsize\bfseries\color{primary}}
  {\thesubsection}{0.5em}{}

\titlespacing{\section}{0pt}{10pt}{4pt}
\titlespacing{\subsection}{0pt}{6pt}{2pt}

% ── Header / Footer ──────────────────────────────────────────
\pagestyle{fancy}
\fancyhf{}
\renewcommand{\headrulewidth}{0pt}
\fancyhead[L]{\small\color{midgray} Data Storm 7.0 — Storming Round}
\fancyhead[R]{\small\color{midgray} Team: \textbf{BigBug}}
\fancyfoot[C]{\small\color{midgray} \thepage}

% ── Hyperlink setup ──────────────────────────────────────────
\hypersetup{
  colorlinks=true,
  linkcolor=primary,
  urlcolor=accent,
  citecolor=primary
}

% ── Helper: info box ─────────────────────────────────────────
\tcbuselibrary{skins, breakable}
\newtcolorbox{infobox}[1][]{
  colback=lightgray,
  colframe=primary,
  fonttitle=\bfseries\small,
  title=#1,
  left=4pt, right=4pt, top=4pt, bottom=4pt,
  boxrule=0.6pt
}

% ── Helper: metric pill table ────────────────────────────────
\newcolumntype{C}[1]{>{\centering\arraybackslash}p{#1}}

% ============================================================
\begin{document}
\thispagestyle{empty}

% ============================================================
% PAGE 1 — COVER PAGE
% ============================================================

\begin{center}
  \vspace*{0.4cm}

  % ── Competition badge ──
  {\color{primary}\rule{\linewidth}{2.5pt}}
  \vspace{0.3cm}

  {\fontsize{9}{11}\selectfont\color{midgray}\MakeUppercase{Sri Lanka's Premier Advanced Analytics Competition}}\\[4pt]
  {\fontsize{22}{26}\selectfont\bfseries\color{primary} Data Storm 7.0}\\[6pt]
  {\fontsize{13}{15}\selectfont\color{accent}\bfseries Powered by OCTAVE JKH}\\[4pt]
  {\color{primary}\rule{\linewidth}{2.5pt}}

  \vspace{1.0cm}

  % ── Report title ──
  {\fontsize{16}{20}\selectfont\bfseries\color{primary}
    Latent Demand Estimation Framework\\[4pt]
    for Traditional Trade Outlets}

  \vspace{0.6cm}
  {\color{midgray}\rule{0.5\linewidth}{0.4pt}}
  \vspace{0.6cm}

  % ── Team details table ──
  \begin{tabular}{rl}
    \textbf{Team Name}   & BigBug \\[3pt]
    \textbf{Members}     & Sithum Fernando, Ranuja Jayawardena, Sandaru Nethmina \\[3pt]
    \textbf{Round}       & Storming Round \\[3pt]
    \textbf{Submission}  & May 17, 2026 \\[3pt]
    \textbf{GitHub}      & \url{https://github.com/The-DarK-Knight-2K/BigBugDataStorm} \\
  \end{tabular}

  \vspace{1.0cm}

  % ── Evaluation weight reminder (visual) ──
  \begin{infobox}[Evaluation Criteria Overview]
    \begin{tabular}{lcc}
      \toprule
      \textbf{Criteria} & \textbf{Weight} & \textbf{Report Section} \\
      \midrule
      Data Engineering \& Forensics & 40\% & Section 1 \\
      Methodology \& Base Math      & 40\% & Sections 2 \& 3 \\
      Generative AI Utilisation     & 20\% & Section 4 \\
      \bottomrule
    \end{tabular}
  \end{infobox}

  \vspace{0.6cm}

  % ── One-line abstract ──
  \begin{minipage}{0.88\linewidth}
    \small\color{midgray}\itshape
    This report presents our end-to-end analytical framework for estimating the
    latent maximum monthly purchase potential (in litres) for 20{,}000 traditional
    trade outlets across Sri Lanka, targeting January 2026.
  \end{minipage}

\end{center}

\newpage

% ============================================================
% PAGE 2 — SECTION 1: DATA FORENSICS & HYGIENE
% ============================================================

\section{Data Engineering \& Forensics}

\subsection{Lakehouse Pipeline Architecture}

We structured our codebase into three explicit layers stored as local
\texttt{.parquet} files under \texttt{Data/Bronze/}, \texttt{Data/Silver/}, and
\texttt{Data/Gold/}.

\begin{center}
\small
\begin{tabular}{p{2.2cm} p{4.5cm} p{5.5cm}}
  \toprule
  \textbf{Layer} & \textbf{Input} & \textbf{Action} \\
  \midrule
  \textcolor{accent}{\textbf{Bronze}} &
    Raw CSV exports from SFA / ERP &
    Ingest as-is; schema-typed; no transformations \\
  \textcolor{primary}{\textbf{Silver}} &
    Bronze tables &
    DQ checks applied; failed records quarantined \\
  \textcolor{successgreen}{\textbf{Gold}} &
    Silver + external POI &
    Feature engineering; model-ready output \\
  \bottomrule
\end{tabular}
\end{center}

\subsection{Reusable Data Quality Check Functions}

All checks are implemented as reusable Python functions in
\texttt{pipeline/Silver/dq\_checks.py} and executed from the Silver pipeline with a
single \texttt{run\_checks()} call per dataset.

\begin{infobox}[DQ Check Summary — Key Functions]
\begin{tabularx}{\linewidth}{l p{5.5cm} X}
  \toprule
  \textbf{Check} & \textbf{Function} & \textbf{Parameters} \\
  \midrule
  Duplicate       & \texttt{duplicate\_check(df, pk\_cols)}
                  & Configurable composite primary key list \\
  Null / Empty    & \texttt{null\_check(df, mandatory\_cols)}
                  & List of non-nullable column names \\
  Referential     & \texttt{ref\_integrity\_check(df, fk\_col, ref\_df)}
                  & Child FK column, parent lookup table \\
  Range           & \texttt{range\_check(df, col, min\_val, max\_val)}
                  & Numeric field + expected bounds \\
  Format / Type   & \texttt{value\_set\_check(df, col, valid\_values)}
                  & Allowed set enforcement for categorical values \\
  \bottomrule
\end{tabularx}
\end{infobox}

\subsection{Anomalies Identified \& Quarantine Results}

\begin{center}
\small
\begin{tabular}{lcccc}
  \toprule
  \textbf{Dataset} & \textbf{Raw Records} & \textbf{Quarantined} & \textbf{Clean} & \textbf{Primary Failure} \\
  \midrule
  transactions\_history & 2,376,389 & 16,620 & 2,359,769 & Negative/invalid volumes \\
  outlet\_master        & 20,000 & 0 & 20,000 & Typos and missing outlet metadata \\
  outlet\_geolocation   & 20,000 & 40 & 19,960 & Invalid / zero GPS coordinates \\
  dist\_seasonality     & 360 & 0 & 370 & Missing Jan 2026 seasonality index \\
  holiday\_list         & 349 & 0 & 351 & Duplicate/multi-type date rows \\
  \bottomrule
\end{tabular}
\end{center}

\textbf{Key System Artefacts Detected:}
\begin{itemize}[noitemsep, topsep=2pt]
  \item \textbf{Ghost entries:} Zero-volume months bounded by active sales,
        suggesting stockouts or system heartbeat records rather than real demand.
  \item \textbf{Clipped volumes:} Repeated identical transaction values were flagged
        as capacity-bound behaviour rather than true consumption signals.
  \item \textbf{Coordinate drift:} 40 outlets had invalid or zero coordinates within
        the raw GPS supply chain and were quarantined before POI enrichment.
  \item \textbf{Master-data decay:} Orphan transaction rows failing referential
        integrity were quarantined to ensure the model only uses validated outlet IDs.
\end{itemize}

\newpage

% ============================================================
% PAGE 3 — SECTION 2: POI DATA ACQUISITION
% ============================================================

\section{POI Data Acquisition}

\subsection{Geospatial Scraping Strategy}

We sourced geospatial Point of Interest (POI) data using the
\textbf{OpenStreetMap Overpass API}. To avoid 20,000 direct point queries, we
engineered a highly scalable, two-phase scraping pipeline.

\textbf{Two-Phase Implementation:}
\begin{enumerate}[noitemsep, topsep=2pt]
  \item \textbf{Phase 1 (Network):} Applied \textbf{K-Means spatial clustering} to group
        20,000 outlets into 400 geographic neighborhoods. We issued one
        bounding-box query (with a 2 km buffer) per cluster and cached the raw
        JSON envelopes locally. This reduced API calls by 98\% and enabled
        safe resumption after failures.
  \item \textbf{Phase 2 (Compute):} A local script parsed the cached JSON and
        calculated precise \textbf{geodesic distances} from each outlet to every
        nearby POI using \texttt{geopy}, correctly accounting for the Earth's
        curvature.
\end{enumerate}

\subsection{POI Categories \& Footfall Rationale}

We extracted 6 core categories. Counts were binned into \textbf{500 m, 1000 m,
and 2000 m} radius bands. The 500 m counts were heavily weighted to generate a
normalized \textbf{Footfall Score (0--100)}.

\begin{center}
\small
\begin{tabularx}{\linewidth}{l l X c}
  \toprule
  \textbf{Category} & \textbf{Core OSM Tags} & \textbf{Demand Signal Rationale} & \textbf{Weight} \\
  \midrule
  Transit     & \texttt{highway=bus\_stop/station} & Constant daily commuter volume & 3.0 \\
  Education   & \texttt{amenity=school/university} & Morning/afternoon density spikes & 2.5 \\
  Markets     & \texttt{shop=supermarket}          & High intent-to-purchase traffic & 2.0 \\
  Healthcare  & \texttt{amenity=hospital/clinic}   & Steady staff \& visitor demand & 1.5 \\
  Community   & \texttt{amenity=place\_of\_worship}& Event-driven congregation footfall & 1.0 \\
  Hospitality & \texttt{tourism=hotel/restaurant}  & Dining crowds and tourists & 1.0 \\
  \bottomrule
\end{tabularx}
\end{center}

\subsection{Final Engineered Features}

Each outlet received 20 newly engineered geospatial features, significantly
upgrading the baseline dataset:
\begin{itemize}[noitemsep, topsep=2pt]
  \item \textbf{18 Dynamic Count Columns:} Specific POI counts for all 6 categories
        across 3 radii (e.g., \texttt{schools\_500m}, \texttt{transport\_1000m},
        \texttt{markets\_2000m}).
  \item \textbf{footfall\_score:} A Min-Max normalized composite score (0--100)
        combining the weighted 500 m metrics.
  \item \textbf{poi\_data\_available:} A boolean data contract flag gracefully
        handling the 40 quarantined outlets lacking valid GPS coordinates by
        assigning them safe zero values.
\end{itemize}

\newpage

% ============================================================
% PAGE 4 — SECTION 3: CAUSAL BASE LOGIC & UNCAPPED DEMAND
% ============================================================

\section{Causal Base Logic \& Uncapped Demand Estimation}

\subsection{Problem Conceptualisation}

Observed historical sales $y_i$ are a \emph{censored} realisation of true
latent demand $y_i^*$:

\begin{equation}
  y_i = \min\!\left(y_i^*,\; C_i\right)
\end{equation}

where $C_i$ represents the binding constraint for outlet $i$ (credit limit,
delivery cap, or stockout ceiling). The goal is to estimate $\hat{y}_i^*$ — the
unconstrained ceiling.

\subsection{Methodology: ``Clean Train, Predict All'' Demand Ceiling}

We implemented a \textbf{``Clean Train, Predict All''} strategy, selecting a gradient-boosted tree model trained on a pseudo-labelled upper-tail
target rather than directly using raw monthly volume. This approximates latent
potential while preserving the signal from constraint-bound outlets.

\begin{infobox}[Model Approach]
  The target is built from the outlet's historical 90th percentile monthly volume
  adjusted for January 2026 seasonality and trading days. We then train a
  CatBoost regressor using 41 Gold-layer structural features. CatBoost was chosen
  after cross-validation demonstrated robust performance (CV RMSE $5.50 \pm 0.38$).
\end{infobox}

\subsection{Feature Set for Potential Model}

\begin{center}
\small
\begin{tabular}{ll}
  \toprule
  \textbf{Feature} & \textbf{Source} \\
  \midrule
  Historical volume signals (p90, max, jan\_avg, ema\_3m) & transactions\_history \\
  Demand volatility (hist\_cv) & transactions\_history \\
  Outlet type, size, and province & outlet\_master \\
  Distributor seasonality multiplier (Jan 2026) & distributor\_seasonality \\
  Activity / Recency (active\_months, months\_since\_last\_order) & transactions\_history \\
  POI counts by category (500m, 1000m, 2000m) & Overpass API \\
  Composite Footfall Score & Overpass API \\
  \bottomrule
\end{tabular}
\end{center}

\subsection{Constraint Neutralisation Logic}

Outlets flagged as constraint-bound (e.g., identical volumes across repeated
periods or zero-volume blackout sequences) had their effective target derived as:

\begin{equation}
  \hat{y}_i^* = \texttt{hist\_p90\_monthly} \times
                \texttt{seasonality\_multiplier\_jan\_2026} \times
                \frac{\texttt{jan\_2026\_trading\_days}}{22.0}
\end{equation}

This uses a robust upper-tail historical signal and normalizes it to the target
month, preserving demand potential even when observed volumes were likely capped.

\newpage

% ============================================================
% PAGE 5 — SECTION 4: GENAI TRANSPARENCY LOG
% ============================================================

\section{Generative AI Transparency Log}

All AI tool usage is documented below in accordance with competition requirements.
Every AI-generated output was critically reviewed and validated before integration.

\begin{center}
\small
\begin{tabularx}{\linewidth}{p{2.8cm} p{1.6cm} X p{2.8cm}}
  \toprule
  \textbf{Task} & \textbf{Tool} & \textbf{How It Was Used} & \textbf{How Validated} \\
  \midrule
  Latent demand framework design &
    Claude &
    Brainstormed pseudo-label vs. censored regression formulations &
    Cross-checked against econometrics literature \\

  Overpass API scraping boilerplate &
    ChatGPT &
    Generated initial Python query builder for OSM POI extraction &
    Manually tested on 10 clusters; fixed bounding-box edge cases \\

  DQ check function templates &
    Copilot &
    Autocompleted parameterisable \texttt{check\_*} functions &
    Unit-tested with synthetic dirty data \\

  Anomaly pattern identification &
    Claude &
    Prompted with sample transaction rows to identify ghost-entry patterns &
    Confirmed via manual inspection of flagged records \\

  Report structure \& LaTeX template &
    Claude &
    Generated section outline and LaTeX skeleton &
    All content written and verified by team \\
  \bottomrule
\end{tabularx}
\end{center}

\subsection*{AI Usage Principles}

\begin{itemize}[noitemsep, topsep=4pt]
  \item \textbf{Accelerator, not a crutch:} AI was used to generate starting
        points and boilerplate; all business logic and model decisions were made
        by the team.
  \item \textbf{Critical evaluation:} Every AI output was tested against real
        data before being incorporated into the pipeline.
  \item \textbf{Iterative prompting:} Prompts were refined multiple times until
        outputs met our quality standards — initial outputs were rarely used
        verbatim.
  \item \textbf{Transparency:} We did not use AI to fabricate results, invent
        data sources, or generate prediction values directly.
\end{itemize}

% ── Footer note ──
\vfill
\begin{center}
  \color{midgray}\small
  \rule{0.4\linewidth}{0.4pt}\\[4pt]
  Submitted for Data Storm 7.0 Storming Round $\cdot$
  Powered by OCTAVE JKH $\cdot$
  Team: \textbf{BigBug}
\end{center}

\end{document}